% - Proposta: é sua arquitetura, descrição, etc...
% O que é ser parametrizável? Muitas vezes alguém faz um projeto de descrição de hardware em vhdl para uma NoC 4x4 de forma rígida. Ou seja, se você quiser sintetizar uma NoC 8x8, você vai ter que editar e refazer vários dos vhdls do projeto. É literalmente um reprojeto. 
% Se a descrição fosse parametrizável, haveria um arquivo de variáveis, onde por exemplo, teria uma que diria o tamanho da NoC=4. Você iria lá, alteraria para 8, e não faria mais nada. Ao fazer a nova síntese, o hardware seria 8x8. Ou seja, o projeto foi feito com uso de variáveis/parâmetros para que uma nova síntese não implique em novo projeto de hardware.

\chapter{The Cart-OnePiece Framework Overview} \label{ch:framework}
\vspace{-1.9cm}

\section{System Architecture and Perception Pipeline} \label{sec:system_pipeline}

The Cart-OnePiece framework is structured as a comprehensive environment designed for synthetic dataset generation, Advanced Autonomous Driving (ADV) workload characterization, and the evaluation of custom task scheduling policies. To achieve these goals, the architecture is built around an asynchronous sensor processing pipeline that closely couples the CARLA simulator \cite{Dosovitskiy2017} with the StarPU runtime system and the NVIDIA TensorRT inference engine. This integration allows the extraction of high-frequency sensor data, its computation across heterogeneous resources, and the observation of the system's performance under distinct scheduling constraints.

A high-level overview of the framework's operational pipeline is presented in Figure~\ref{fig:framework_pipeline}. The system operates on a client-server model where the CARLA server acts as the primary workload generator, simulating the physics, traffic, and environmental conditions of the urban scenario while rendering the virtual world from the perspective of the ego-vehicle. On the client side, a C++ application implemented on top of the StarPU library dictates the asynchronous extraction and processing of these sensor streams.

\begin{figure}[hbt!]
  \centering
  \includegraphics[width=\textwidth]{Figure2.png}
  \caption{Overview of the Cart-OnePiece framework and the autonomous driving perception pipeline.}
  \label{fig:framework_pipeline}
\end{figure}

The core of the execution flows through a continuous perception pipeline composed of consecutive computational stages. In the initial phase, multiple virtual cameras attached to the simulated vehicle output raw visual data (typically in RGBA format) and their corresponding semantic segmentation ground truth. Rather than blocking the main execution thread to process these images immediately, the StarPU runtime is tasked with managing these operations asynchronously, mapping them dynamically to the available CPU workers and GPU resources.

Once the raw data is captured, a preprocessing stage converts the images from the simulator's native BGRA format into the NCHW tensor layout required by deep learning engines, normalizing the pixel values simultaneously. This stage is highly parallelizable and is primarily distributed among the host CPU cores. Following the preprocessing, the normalized tensors are dispatched to the inference stage. Here, the TensorRT library, accelerated via CUDA, processes the data using pre-compiled models, such as ResNet50 or DeepLabV3, predicting the semantic masks of the scene.

The resulting semantic predictions are then evaluated against the synchronously captured ground truth in a metric calculation stage, which computes the Mean Intersection over Union (mIoU) and pixel accuracy. To ensure that the computational trace is accurately registered without interfering with the workload execution, the pipeline relies on the APEX profiling tool and StarPU's internal FxT tracing mechanisms, which record the latency, memory utilization, and structural task dependencies of each frame.

This architectural design not only provides a flexible testbed for characterizing the ADV workload but also exposes a mechanism to control the depth of in-flight tasks and constrain resource mappings. By keeping the different processing stages decoupled through the StarPU runtime, the Cart-OnePiece framework enables a direct evaluation of how different scheduling policies impact the end-to-end latency and the efficiency of heterogeneous component utilization.


